\documentclass{rapport}
\usepackage{lipsum}
\usepackage{gensymb}
\usepackage{float}
\usepackage{graphicx} % Required for inserting images
\usepackage{listings}
\usepackage{xcolor}
\usepackage{minted}
\usepackage{amsmath, amssymb, amsfonts}
\usepackage{hyperref}

\pdfminorversion=5
\pdfcompresslevel=9
\pdfobjcompresslevel=5

\hypersetup{
    colorlinks=true,        % Enable colored links
    linkcolor=blue,         % Color for internal links
    citecolor=black,         % Color for citations
    filecolor=blue,         % Color for file links
    urlcolor=blue           % Color for URLs
}

\title{file title} %title of the file

\definecolor{codegreen}{rgb}{0,0.6,0}
\definecolor{codegray}{rgb}{0.5,0.5,0.5}
\definecolor{codepurple}{rgb}{0.58,0,0.82}
\definecolor{backcolour}{rgb}{0.95,0.95,0.92}

\lstdefinestyle{mystyle}{
    backgroundcolor=\color{backcolour},   
    commentstyle=\color{codegreen},
    keywordstyle=\color{magenta},
    numberstyle=\tiny\color{codegray},
    stringstyle=\color{codepurple},
    basicstyle=\ttfamily\footnotesize,
    breakatwhitespace=false,         
    breaklines=true,                 
    captionpos=b,                    
    keepspaces=true,                 
    numbers=left,                    
    numbersep=5pt,                  
    showspaces=false,                
    showstringspaces=false,
    showtabs=false,                  
    tabsize=2
}

\lstset{style=mystyle}

\begin{document}

%----------- Report information ---------
\logo{logos/uoft_logo.png}
\uni{\textbf{University of Toronto}}
\ttitle{AER1515 Perception for Robotics: Assignment 3} %title of the file
\subject{AER1515} % Subject name
\topic{Assignment 2} % Topic name

\students{\textsc{Yang-Chen Lin}} % information related to the students

%----------- Init -------------------
        
\buildmargins % display margins
\buildcover % create the front cover of the document
\toc % creates the table of contents

%------------ Report body ----------------
\section{Answers}
\subsection{Part 1: Depth Estimation}
\begin{figure}[H]
    \centering
    \includegraphics[width=0.63\textwidth]{images/est_depth_stack.png}
    \caption{Estimateded depth using dense disparity maps. The color bar represents the depth in meters.}
    \label{fig:1}
\end{figure}
\paragraph{Observations on the Depth Image Quality}
The produced depth images show reasonable consistency across the scene, with smooth gradients along the road and clear separation between 
near and far objects. Depth accuracy is generally higher for textured and well-lit surfaces, such as the road and nearby vehicles, 
where disparity estimation remains stable and consistent. In contrast, reflective or low-texture areas, tree canopies, and occlusion 
boundaries often exhibit noisy or discontinuous depth values due to unreliable or missing disparity information.

\subsection{Part 2: 2D Bounding Box Detection}
\begin{figure}[H]
    \centering
    \includegraphics[width=0.68\textwidth]{images/yolo_detections_stack.png}
    \caption{Estimated 2D bounding box using the YOLOv3 object detector with \texttt{confidence\_threshold:0.51} and \texttt{NMS\_threshold:0.6}}
    \label{fig:2}
\end{figure}

\subsection{Part 3: Instance Segmentation}
To perform the instance segmentation task, I implemented two following approaches:
\begin{enumerate}
    \item Basic Method
    \item YOLOv11x-seg model
\end{enumerate}

\subsubsection{Basic Method}
A typical non-learning method to perform instance segmentation is to:
\begin{enumerate}
    \item Calculate the average depth within the given bounding boxes
    \item Set a distance threshold from the computed average depth and assign fitting pixels as the object's segmentation mask 
\end{enumerate}
The result are displayed below. Note that for the python scripts to generate segmentation masks using the basic methods 
are limited to the train datasets, as the performance are deemed to be worse than 
learning methods, so I utilized the ground truth depth map and bounding boxes provided with the data.
\begin{figure}[H]
    \centering
    \begin{minipage}{0.48\textwidth}
        \centering
        \includegraphics[width=\textwidth]{images/yolo_detections_basic.png} 
        \subcaption{2D bounding boxes for train images.}
    \end{minipage}
    \begin{minipage}{0.48\textwidth}
        \centering
        \includegraphics[width=\textwidth]{images/est_seg_basic_stack.png}
        \subcaption{Basic instance segmentation method with \texttt{distance\_threshold:5.8}}
    \end{minipage}
    \caption{}
    \label{fig:3}
\end{figure}

From \href{fig:3}{Fig. 3}, we can see that the segmentation mask genreated by the basic method aren't very accurate. Therefore, it's
time to utilize the power of deep learning methods.

\subsubsection{YOLOv11x-seg Instance Segmentation Model}
YOLOv11x-seg is an advanced instance segmentation model developed by Ultralytics \cite{yolo11_ultralytics}, based on the YOLO (You Only Look Once) 
architecture. It performs real-time object detection and segmentation in a single forward pass through the network. The model 
consists of a convolutional backbone for feature extraction, a neck with a feature pyramid network (FPN) for multi-scale feature 
fusion, and a detection head that predicts bounding boxes, class probabilities, and pixel-level masks for each detected object. 
Trained on large datasets like COCO, it uses anchor-free predictions and advanced loss functions to achieve high accuracy and 
speed. The following \href{fig:4}{Fig. 4} shows the 2D bounding box and instance segmentation masks using this architecture.

\begin{figure}[H]
    \centering
    \begin{minipage}{0.48\textwidth}
        \centering
        \includegraphics[width=\textwidth]{images/est_bbox_yolov11_stack.png} 
        \subcaption{2D bounding boxes detected by the YOLOv11x model.}
    \end{minipage}
    \begin{minipage}{0.48\textwidth}
        \centering
        \includegraphics[width=\textwidth]{images/est_seg_yolov11x_stack.png}
        \subcaption{Instance segmentation masks for cars; generated by the YOLOv11x model.}
    \end{minipage}
    \caption{}
    \label{fig:4}
\end{figure}
\bibliographystyle{ieeetr}  
\bibliography{references} 

\end{document}
